\section{Schroedinger-, Heisenberg- und Dirac/Wechselwirkungs-Bild }

In der Quantenmechanik kann dieselbe Physik in verschiedenen Darstellungen formuliert werden. Die am haeufigsten verwendeten sind das Schroedinger-Bild, das Heisenberg-Bild und das Dirac-Bild (Interaction Picture). Sie unterscheiden sich nur darin, welcher Anteil der Zeitentwicklung den Zustaenden bzw. den Operatoren zugeordnet wird. Messbare Groessen sind in allen drei Bildern identisch.

\subsubsection*{Schroedinger-Bild}

Im Schroedinger-Bild tragen die Zustaende die Zeitabhaengigkeit, waehrend Operatoren (sofern sie nicht explizit zeitabhaengig sind) zeitunabhaengig bleiben. Fuer einen Zustand $\ket{\psi_S(t)}$ gilt
\begin{equation}
    i\hbar\,\frac{\partial}{\partial t}\ket{\psi_S(t)}=H\ket{\psi_S(t)}.
\end{equation}
Die Zeitableitung eines Operators ergibt sich nur aus seiner expliziten Zeitabhaengigkeit:
\begin{equation}
    \frac{d}{dt}A_S=\frac{\partial}{\partial t}A_S.
\end{equation}
Ist der Hamiltonoperator zeitunabhaengig, dann ist die Loesung
\begin{equation}
    \ket{\psi_S(t)}=U(t,t_0)\ket{\psi_S(t_0)},
    \qquad
    U(t,t_0)=e^{-\frac{i}{\hbar}H(t-t_0)}.
\end{equation}
Fuer Erwartungswerte folgt
\begin{equation}
    \langle A\rangle_t=\bra{\psi_S(t)}A_S\ket{\psi_S(t)}.
\end{equation}
Setzt man $\ket{\psi_S(t)}=U(t,t_0)\ket{\psi_S(t_0)}$ ein, erhaelt man die direkte Gegenueberstellung von Schroedinger- und Heisenberg-Bild:
\begin{equation}
    \begin{aligned}
        \langle A\rangle_t
         & = \underbrace{\bra{\psi_S(t_0)}U^{\dagger}(t,t_0)}_{\text{Schroedinger: Zustand}}
        \,\underbrace{A_S}_{\text{Operator}}\,
        \underbrace{U(t,t_0)\ket{\psi_S(t_0)}}_{\text{Schroedinger: Zustand}}                \\
         & = \underbrace{\bra{\psi_H}}_{\text{Heisenberg: Zustand}}
        \,\underbrace{\bigl(U^{\dagger}(t,t_0)A_SU(t,t_0)\bigr)}_{\text{Heisenberg: Operator }A_H(t)}\,
        \underbrace{\ket{\psi_H}}_{\text{Heisenberg: Zustand}}.
    \end{aligned}
\end{equation}

\subsubsection*{Heisenberg-Bild}

Im Heisenberg-Bild wird die Zeitentwicklung vollstaendig auf die Operatoren verschoben. Die Zustaende sind (bei zeitunabhaengigem $H$) konstant, z.B.
\begin{equation}
    \ket{\psi_H}=\ket{\psi_S(t_0)}.
\end{equation}
Damit gilt unmittelbar
\begin{equation}
    \frac{d}{dt}\ket{\psi_H}=0.
\end{equation}
Die zugehoerigen Operatoren lauten
\begin{equation}
    A_H(t)=U^{\dagger}(t,t_0)A_SU(t,t_0).
\end{equation}
Fuer zeitunabhaengiges $H$ ist das aequivalent zu
\begin{equation}
    A_H(t)=e^{\frac{i}{\hbar}H(t-t_0)}A_Se^{-\frac{i}{\hbar}H(t-t_0)}.
\end{equation}
Daraus folgt die Heisenberg-Gleichung
\begin{equation}
    i\hbar\,\frac{d}{dt}A_H(t)=\left[A_H(t),H\right]+i\hbar\left(\frac{\partial A_S}{\partial t}\right)_H.
\end{equation}
Fuer Operatoren ohne explizite Zeitabhaengigkeit bleibt nur der Kommutatorterm. Erwartungswerte sind unveraendert:
\begin{equation}
    \langle A\rangle_t=\bra{\psi_H}A_H(t)\ket{\psi_H}.
\end{equation}

\subsubsection*{Dirac-Bild / Interaction Picture}

Das Interaction Picture ist besonders nuetzlich fuer Stoerungstheorie und Feynman-Diagramme. Man zerlegt den Hamiltonoperator in
\begin{equation}
    H=H_0+V,
\end{equation}
wobei $H_0$ der exakt loesbare Teil und $V$ die Wechselwirkung ist. Dann wird die "freie" Dynamik mit $H_0$ auf die Operatoren gelegt, waehrend die zusaetzliche Dynamik durch $V$ in den Zustaenden bleibt:
\begin{equation}
    A_I(t)=e^{\frac{i}{\hbar}H_0(t-t_0)}A_Se^{-\frac{i}{\hbar}H_0(t-t_0)},
\end{equation}
\begin{equation}
    \ket{\psi_I(t)}=e^{\frac{i}{\hbar}H_0(t-t_0)}\ket{\psi_S(t)}.
\end{equation}
Die Zustaende erfuellen damit
\begin{equation}
    i\hbar\,\frac{\partial}{\partial t}\ket{\psi_I(t)}=V_I(t)\ket{\psi_I(t)},
\end{equation}
mit
\begin{equation}
    V_I(t)=e^{\frac{i}{\hbar}H_0(t-t_0)}Ve^{-\frac{i}{\hbar}H_0(t-t_0)}.
\end{equation}
Die zugehoerige Zeitentwicklungsoperatorgleichung lautet
\begin{equation}
    i\hbar\,\frac{\partial}{\partial t}U_I(t,t_0)=V_I(t)U_I(t,t_0),
    \qquad U_I(t_0,t_0)=\mathds{1}.
\end{equation}
Ihre formale Loesung ist die Dyson-Reihe
\begin{equation}
    U_I(t,t_0)=\mathcal{T}\exp\left[-\frac{i}{\hbar}\int_{t_0}^{t}dt'\,V_I(t')\right],
\end{equation}
mit dem Zeitordnungsoperator $\mathcal{T}$. Genau diese Struktur wird in der Vielteilchenphysik fuer Greensche Funktionen und Diagrammtechniken verwendet.

\subsubsection*{Physikalische Einordnung}

Die drei Bilder sind mathematisch aequivalente Beschreibungen derselben Theorie. Das Schroedinger-Bild ist oft der intuitivste Startpunkt. Das Heisenberg-Bild ist fuer Operatorgleichungen und Erhaltungssaetze besonders klar. Das Interaction Picture kombiniert beide Vorteile: freie Ausbreitung steckt in den Operatoren, Wechselwirkungen in den Zustaenden. Deshalb ist es die natuerliche Sprache fuer Streuprozesse, Selbstenergien, renormierte Quasiteilchen und die systematische Entwicklung in Stoerungsordnungen.
